% Standalone resolution manuscript generated from solution.md.
% Compile with XeLaTeX; author and review metadata are maintained in Markdown.
\documentclass[10pt,a4paper]{article}
\usepackage[margin=24mm,headheight=15pt]{geometry}
\usepackage{fontspec}
\setmainfont{DejaVuSerif.ttf}[BoldFont=DejaVuSerif-Bold.ttf,ItalicFont=DejaVuSerif-Italic.ttf,BoldItalicFont=DejaVuSerif-BoldItalic.ttf]
\setsansfont{DejaVuSans.ttf}[BoldFont=DejaVuSans-Bold.ttf,ItalicFont=DejaVuSans-Oblique.ttf]
\setmonofont{DejaVuSansMono.ttf}
\usepackage{amsmath,amssymb,unicode-math}
\setmathfont{latinmodern-math.otf}
\usepackage{xcolor,microtype,parskip,needspace,fancyhdr}
\usepackage{bookmark}
\usepackage{xurl}
\hypersetup{colorlinks=true,urlcolor=blue,linkcolor=blue,pdftitle={IE-15:
exact rook-pivoting growth factors in orders three and
four},pdfauthor={George Stepaniants},pdflang={en-GB}}
\urlstyle{same}
\setlength{\emergencystretch}{3em}
\providecommand{\tightlist}{\setlength{\itemsep}{0pt}\setlength{\parskip}{0pt}}
\providecommand{\pandocbounded}[1]{#1}
\setcounter{secnumdepth}{0}
\allowdisplaybreaks[1]
\widowpenalty=10000
\clubpenalty=10000
\pagestyle{fancy}
\fancyhf{}
\fancyhead[L]{\small\sffamily Verified resolution / IE-15}
\fancyhead[R]{\small\sffamily 11 September 2026}
\fancyfoot[L]{\parbox[b]{0.9\textwidth}{\fontsize{8}{10}\selectfont Independent
Codex-agent review; not external human peer review or formal
certification.}}
\fancyfoot[R]{\footnotesize\thepage}
\begin{document}
{\raggedright\Large\bfseries IE-15: exact rook-pivoting growth factors
in orders three and four\par}
\vspace{0.4em}
{\normalsize\bfseries George Stepaniants\par}
{\small Department of Computing and Mathematical Sciences, California
Institute of Technology, Pasadena, California, USA\par}
\vspace{0.6em}
\pagestyle{plain}

\textbf{Verification status: independent-agent PASS.} Developed with
ChatGPT/Codex, with numerical searches used for discovery and an
analytic proof of the universal bounds. A separate Codex agent audited
the complete proof and independently checked the exact witnesses. This
is not external human peer review or formal certification.
\href{https://github.com/ajt60gaibb/OpenProblemsInNLA/blob/main/references/stepaniants-ie15-2026-09-11/verification/reviews/IE-15-review.md}{Detailed
review}.

All matrices are real and all arithmetic is exact. A rook pivot is any
nonzero entry maximal in absolute value in its active row and active
column. All admissible choices, including ties, are allowed, exactly as
in canonical IE-15.

If \(S_1=A,\ldots,S_n\) are the active matrices for a path \(\pi\),
define

\[
\rho(A,\pi)=\frac{\max_{1\le k\le n}\|S_k\|_{\max}}{\|A\|_{\max}},
\qquad \|M\|_{\max}=\max_{i,j}|M_{ij}|,
\]

and let \(g_{\mathrm{RP}}(n)\) be the supremum of \(\rho(A,\pi)\) over
real nonsingular \(A\in\mathbb R^{n\times n}\) and all its admissible
rook paths.

\textbf{Theorem 1 (exact small-order rook-pivoting growth).} For real
nonsingular matrices and every admissible rook-pivoting path, with all
ties allowed and growth measured over every active-matrix entry as in
IE-15,

\[
\boxed{g_{\mathrm{RP}}(3)=3,\qquad g_{\mathrm{RP}}(4)=14/3.}
\]

Both maxima are attained by the rational matrices in Section 5.

\subsection{1. Normalization and elimination
coordinates}\label{normalization-and-elimination-coordinates}

Fix any admissible path on a nonsingular matrix, divide the matrix by
its initial maximum absolute entry, and apply the path's eventual row
and column permutations at the outset. This reproduces the active
matrices up to row and column permutations and lets us take all pivots
in diagonal order. Gaussian elimination then gives

\[
A=LDR,
\]

where \(L\) and \(R\) are unit lower and unit upper triangular,
respectively, and \(D=\operatorname{diag}(p_1,\ldots,p_n)\). The rook
conditions are exactly

\[
|L_{ij}|\le1\quad(i>j),\qquad |R_{ij}|\le1\quad(i<j).
\tag{1}
\]

Indeed, the active pivot column and row at step \(j\) are \(p_j L_{ij}\)
and \(p_jR_{jk}\). All original entries satisfy \(|A_{ij}|\le1\).

Multiplying suitable original rows by \(-1\) makes every \(p_j\)
positive, preserving all absolute-value conditions and growth.
Specifically, if \(S\) is diagonal with entries
\(\operatorname{sign}(p_j)\), then \(SA=(SLS)(SD)R\). Next, a
simultaneous row and column sign change \(A\mapsto QAQ\) preserves the
positive pivots and can arrange

\[
L_{ni}\ge0\qquad(1\le i<n).
\tag{2}
\]

Finally, when bounding the positive last pivot, we may replace
\(A_{nn}\) by \(1\). This leaves every earlier pivot and its active row
and column unchanged, so rook admissibility persists. The final pivot
only increases by \(1-A_{nn}\), and the matrix stays nonsingular. Thus
it suffices to bound the last pivot with \(A_{nn}=1\).

The first pivot obeys \(0<p_1\le1\), and

\[
0<p_2\le1+p_1\le2,
\tag{3}
\]

because \(A_{22}=p_2+p_1L_{21}R_{12}\). Every entry after one
elimination step is at most \(1+p_1\le2\) in absolute value. Every entry
after two steps is at most \(1+p_1+p_2\le4\). These bounds follow
directly by subtracting the corresponding rank-one updates, using (1).

\subsection{2. The order-three upper
bound}\label{the-order-three-upper-bound}

Use the normalization above with \(n=3\), and write

\[
p=p_1,\quad q=p_2,\quad a=L_{21},\quad b=R_{12},
\quad c=L_{31}\ge0,\quad e=L_{32}\ge0,
\quad d=-R_{13},\quad f=-R_{23}.
\]

All six multipliers lie in \([-1,1]\), with \(c,e\in[0,1]\). The last
pivot is

\[
p_3=1+pcd+qef.
\]

Suppose \(p_3>3\). Since \(pcd\le1\) and \(qef\le2\), this implies
\(pcd>0\) and \(qef>1\). Consequently \(d,f>0\) and both \(qe>1\) and
\(qf>1\). But

\[
A_{23}=-qf-pad,\qquad A_{32}=qe+pcb.
\]

The bounds \(|A_{23}|,|A_{32}|\le1\) force \(a<0\) and \(b<0\),
respectively. This contradicts

\[
A_{22}=q+pab>q>1.
\]

Thus \(p_3\le3\). Intermediate entries are bounded by two, so
\(g_{\mathrm{RP}}(3)\le3\).

The same bound on the final two-step Schur value holds for a possibly
singular order-three matrix provided its first two pivots are nonzero
and admissible. If that final value is nonzero, the argument above
applies after a last-row sign change; if it is zero, the bound is
immediate. We will use this observation for an order-three submatrix in
Section 4.

\Needspace{20\baselineskip}

\subsection{3. A two-pivot scalar
inequality}\label{a-two-pivot-scalar-inequality}

Define

\[
h(c,d)=3cd+1+|c-d|.
\]

\textbf{Lemma.} Suppose

\[
0<p\le1,\quad q>0,\quad a,b\in[-1,1],
\quad c_1,c_2\in[0,1],\quad d_1,d_2\in[-1,1],
\]

and

\[
|q+pab|\le1,\qquad
|qd_2+pa d_1|\le1,\qquad
|qc_2+pc_1b|\le1.
\tag{4}
\]

Then

\[
p\,h(c_1,d_1)+q\,h(c_2,d_2)\le8.
\tag{5}
\]

\textbf{Proof.} We use the following elementary inequalities on the
indicated domains:

\[
0\le h(c,d)\le4,\qquad h(c,d)\le2+2c
\quad(0\le c\le1,\ -1\le d\le1),
\tag{6}
\]

\[
h(c,d)\le2+2d\quad(d\ge0),\qquad
h(c,d)\le2\quad(d\le0),
\tag{7}
\]

\[
h(c,d)\le2+2cd\quad(0\le c,d\le1).
\tag{8}
\]

For completeness, convexity in either variable gives the upper bounds in
(6) and (7) by evaluating the endpoints. When \(d=-v\le0\),
\(h(c,d)=1+c+v-3cv\) lies between zero and two by the four corners of
the unit square. For \(d\ge0\), nonnegativity is immediate. Inequality
(8) follows from \(|c-d|\le1-cd\).

If \(q\le1\), (5) follows from \(h\le4\). Hence assume \(q>1\). The
first condition in (4) implies \(ab<0\) and

\[
q\le1+p|ab|\le1+p\le2.
\tag{9}
\]

If \(b>0\), then \(a<0\) and the last condition in (4) gives
\(qc_2+pbc_1\le1\). Using (6),

\[
\begin{aligned}
p h(c_1,d_1)+q h(c_2,d_2)
&\le2(p+q)+2(pc_1+qc_2)\\
&\le2\bigl(p+q+1+p(1-b)c_1\bigr)\\
&\le2\bigl(p+q+1+p(1-b)\bigr)\\
&\le4+4p\le8,
\end{aligned}
\]

where the penultimate step uses \(q\le1+pb\).

It remains to consider \(a>0\), \(b<0\). If \(d_2\le0\), then (6), (7),
and (9) give

\[
p h(c_1,d_1)+q h(c_2,d_2)
\le4p+2q\le2+6p\le8.
\]

If \(d_2>0\) and \(d_1\ge0\), then \(qd_2+pa d_1\le1\). Using (7) in
place of (6), the same calculation as above gives

\[
p h(c_1,d_1)+q h(c_2,d_2)
\le2\bigl(p+q+1+p(1-a)\bigr)\le4+4p\le8,
\]

because \(q\le1+pa\).

The only remaining case is \(d_1=-v<0<d_2\). Put \(c=c_1\), \(B=-b\),
and

\[
U=1+pBc,\qquad V=1+pav.
\]

The last two conditions in (4) give \(qc_2\le U\) and \(qd_2\le V\).
Since \(0\le c_2,d_2\le1\),

\[
q c_2d_2\le\min\{q,U,V,UV/q\}.
\]

By (8), the left side of (5) is therefore at most

\[
p\,h(c,-v)+\Phi(q,U,V),\qquad
\Phi(q,U,V)=2q+2\min\{q,U,V,UV/q\}.
\]

For positive arguments, \(\Phi\) is nondecreasing in each argument.
Monotonicity in \(U,V\) follows termwise. To check the \(q\) argument,
put \(u=\min(U,V)\), \(w=\max(U,V)\): the formula for \(\Phi\) is \(4q\)
when \(q\le u\), \(2q+2u\) when \(u\le q\le w\), and \(2q+2UV/q\) when
\(q\ge w\). In the last region its derivative is \(2-2UV/q^2\ge0\), and
the three pieces join continuously.

Now \(q\le2\), \(U\le1+c\), \(V\le1+v\), and
\(0\le p\,h(c,-v)\le h(c,-v)\). Hence

\[
\begin{aligned}
p h(c_1,d_1)+q h(c_2,d_2)
&\le h(c,-v)+\Phi(2,1+c,1+v)\\
&=1+c+v-3cv+4+(1+c)(1+v)\\
&=8-2(1-c)(1-v)\le8.
\end{aligned}
\]

This proves every case of the lemma. \(\square\)

\subsection{4. The order-four upper
bound}\label{the-order-four-upper-bound}

Use Section 1 with \(n=4\), positive pivots, \(A_{44}=1\), and

\[
c_i=L_{4i}\in[0,1],\qquad d_i=-R_{i4}\in[-1,1],
\qquad w_i=p_i c_i d_i\quad(i=1,2,3).
\]

Then

\[
p_4=1+w_1+w_2+w_3.
\tag{10}
\]

Write \(W=w_1+w_2\). The principal submatrix on indices \(1,2,4\) has
the same first two pivots and retains their rook admissibility, since
deleting other rows and columns cannot violate those maximum conditions.
Its final Schur value is \(1+W\). The order-three result, including its
possibly singular extension, gives

\[
W\le2.
\tag{11}
\]

If \(d_3\le0\) or \(c_3=0\), then \(w_3\le0\) and \(p_4\le3\). We may
therefore assume \(c_3,d_3\in(0,1]\), and denote these two numbers
temporarily by \(C,D\).

For \(i=1,2\), put \(u_i=L_{3i}\), \(v_i=R_{i3}\). The original-entry
inequalities \(A_{33}\le1\), \(-A_{34}\le1\), and \(A_{43}\le1\) are

\[
p_3+\sum_{i=1}^2p_i u_i v_i\le1,
\qquad p_3D+\sum_{i=1}^2p_i u_i d_i\le1,
\qquad p_3C+\sum_{i=1}^2p_i c_i v_i\le1.
\]

\Needspace{9\baselineskip}

Multiply these inequalities by \(CD,C,D\), respectively, and add. Then

\[
3w_3\le CD+C+D-
\sum_{i=1}^2p_i(CD u_iv_i+C d_i u_i+D c_i v_i).
\]

Since \(|u_i|,|v_i|\le1\), define

\[
\mathcal B(C,D)=CD+C+D+
\sum_{i=1}^2p_i\left[3c_id_i-
\min_{-1\le u,v\le1}(CDuv+C d_i u+D c_i v)\right].
\]

We have

\[
3(W+w_3)\le\mathcal B(C,D).
\tag{12}
\]

The bilinear minimum over \((u,v)\) is attained at the four corners of
their square. Its negative is a maximum of four functions, each affine
in \(C\) when \(D\) is fixed and affine in \(D\) when \(C\) is fixed.
Thus \(\mathcal B\) is convex separately in \(C\) and \(D\). Two
applications of the one-variable endpoint bound show that its maximum on
\([0,1]^2\) is bounded by its maximum at the four corners.

Those corners have the following bounds, using (3) and (11):

\[
\mathcal B(0,0)=3W\le6,
\]

\[
\mathcal B(1,0)=1+3W+\sum_{i=1}^2p_i|d_i|
\le1+6+(p_1+p_2)\le10,
\]

\[
\mathcal B(0,1)=1+3W+\sum_{i=1}^2p_i c_i\le10.
\]

For the remaining corner, the four elementary possibilities give

\[
\min_{-1\le u,v\le1}(uv+du+cv)=-1-|c-d|
\quad(0\le c\le1,\ -1\le d\le1).
\]

Indeed the two opposite-sign corners give the right-hand side; the two
same-sign corners are no smaller. Consequently

\[
\mathcal B(1,1)=3+p_1h(c_1,d_1)+p_2h(c_2,d_2)\le11.
\tag{13}
\]

The last inequality is precisely the lemma, with \(p=p_1\), \(q=p_2\),
\(a=L_{21}\), and \(b=R_{12}\). Its three assumptions (4) are the
original bounds on \(A_{22}\), \(-A_{24}\), and \(A_{42}\),
respectively.

Every corner is therefore at most eleven. Equations (10) and (12) yield

\[
p_4=1+(W+w_3)\le1+\frac{11}{3}=\frac{14}{3}.
\]

Entries at the earlier stages are bounded by one, two, and four, as
shown in Section 1. Since \(4<14/3\), this establishes the full growth
bound \(g_{\mathrm{RP}}(4)\le14/3\), including every intermediate entry
and every admissible choice and tie.

\subsection{5. Matching exact examples}\label{matching-exact-examples}

For order three take

\[
A_3=\begin{pmatrix}1&0&-1\\0&1&-1\\1&1&1\end{pmatrix}.
\]

Its initial maximum absolute entry is one. Diagonal elimination gives

\[
S_2=\begin{pmatrix}1&-1\\1&2\end{pmatrix},\qquad S_3=(3).
\]

The first and second pivots are one and are maximal in their active rows
and columns. Thus the path is admissible and has growth three. Its
determinant is three, so the matrix is nonsingular.

For order four take

\[
A_4=\begin{pmatrix}
1&0&1&1\\
0&1&1/3&-1\\
-1/3&-1&1&-1\\
-1&1&1&1
\end{pmatrix}.
\]

Again the initial maximum absolute entry is one. Its successive active
matrices under the diagonal path are

\[
S_2=\begin{pmatrix}
1&1/3&-1\\
-1&4/3&-2/3\\
1&2&2
\end{pmatrix},\qquad
S_3=\begin{pmatrix}5/3&-5/3\\5/3&3\end{pmatrix},\qquad
S_4=(14/3).
\]

Each selected pivot is maximal in its active row and column; the ties
are expressly allowed by the problem. The pivot sequence is
\(1,1,5/3,14/3\), whose product is \(70/9\ne0\). The growth is exactly
\(14/3\). Together with the universal upper bounds, these witnesses
prove

\[
\boxed{g_{\mathrm{RP}}(3)=3,\qquad g_{\mathrm{RP}}(4)=14/3.}
\]

\subsection{Scope and computational
evidence}\label{scope-and-computational-evidence}

The upper bounds are analytic, not numerical optimization conclusions.
Numerical constrained optimization was used only to discover the
order-four witness and suggest the sharp bound. The proof includes all
real nonsingular matrices, arbitrary admissible rook paths, all ties,
all intermediate active entries, and matching attained lower bounds. No
particular tie-breaking or pivot-search implementation is assumed.

The standard-library Python script
\href{https://github.com/ajt60gaibb/OpenProblemsInNLA/blob/main/references/stepaniants-ie15-2026-09-11/verification/verify_witnesses.py}{\texttt{verify\_witnesses.py}}
checks both rational witnesses: nonsingularity, every selected pivot's
row and column inequalities, all active matrices, and the attained
growth. It uses \texttt{fractions.Fraction} throughout. These finite
checks supplement the analytic upper-bound proof; they are not a
substitute for it.

\subsection{References}\label{references}

\begin{enumerate}
\def\labelenumi{\arabic{enumi}.}
\tightlist
\item
  \href{https://github.com/ajt60gaibb/OpenProblemsInNLA/blob/main/linear-systems-and-elimination/IE-15/README.md}{IE-15
  --- Exact small-order growth factors for rook pivoting}, \emph{Open
  Problems in Numerical Linear Algebra}. This is the exact
  all-admissible-path target settled by Theorem 1.
\item
  N. J. Higham,
  \href{https://doi.org/10.1137/1.9780898718027}{\emph{Accuracy and
  Stability of Numerical Algorithms}, second edition}, SIAM (2002),
  Problem 9.18, p.~193. This is the historical small-order question
  cited by IE-15; no result in the book is used to prove the upper
  bounds here.
\item
  A. Edelman and J. Urschel,
  \href{https://doi.org/10.1137/23M1571903}{\emph{Some New Results on
  the Maximum Growth Factor in Gaussian Elimination}}, SIAM J. Matrix
  Anal. Appl. 45 (2024),
  \href{https://arxiv.org/html/2303.04892v4}{Section 6, primary
  manuscript}. Context for the rook-pivoting model and broader
  growth-factor problem.
\item
  R. Shah and J. Urschel,
  \href{https://arxiv.org/html/2608.19189v4}{\emph{Entry growth in
  Gaussian elimination}}, arXiv:2608.19189v4, August 2026. Context for
  general-order growth; the finite-order proof above is independent of
  its asymptotic results.
\end{enumerate}
\end{document}
